\chapter*{Your Linux Toolkit and Environment Setup}
\addcontentsline{toc}{chapter}{Your Linux Toolkit and Environment Setup}
\markboth{Your Linux Toolkit and Environment Setup}{}
\chapterguide{No Linux experience is required. You need a computer where you can use Linux, WSL, or a Linux virtual machine.}{identify the course folder, open a Linux shell, edit source in Vim, compile with GCC, run a program, and prepare Git/SSH for saving course work.}

The practicals use a \textbf{command-line workflow}. Linux is the reference environment because the processes, files, permissions, IPC objects, and other OS concepts used in the course are easy to observe there.

\section*{The course project layout}
Keep the supplied structure unchanged. The lab pages print source directly from \texttt{all\_experiments/}, so the files in the book and the files you edit are the same files.

\begin{lstlisting}[style=osbash]
OPERATING SYSTEMS/
|-- Lab01/
|   |-- directory_list.c
|   |-- fork_demo.c
|-- Lab02/
|   |-- file_copy.c
|   |-- file_info.c
|   |-- mini_grep.c
|   |-- sample.txt
|-- Lab03/
|   |-- events.log
|   |-- log_summary.sh
|   |-- process_summary.sh
|   |-- system_report.sh
|-- Lab04/
|   |-- fcfs.c
|   |-- priority.c
|   |-- round_robin.c
|   |-- sjf.c
|-- Lab05/
|   |-- pipe_demo.c
|   |-- producer_consumer_sim.c
|   |-- shm_reader.c
|   |-- shm_writer.c
|-- Lab06/
|   |-- banker.c
|   |-- threads_demo.c
|-- Lab07/
|   |-- best_fit.c
|   |-- first_fit.c
|   |-- worst_fit.c
|-- Lab08/
|   |-- fifo.c
|   |-- lfu.c
|   |-- lru.c
|-- Lab09/
|   |-- directory_walk.c
|   |-- make_demo_tree.c
|-- Lab10/
|   |-- binary_records.c
|   |-- indexed_allocation.c
|   |-- text_records.c
|-- README.md
\end{lstlisting}

\section*{Choose the Linux environment}
\begin{mltable}
\renewcommand{\arraystretch}{1.18}
\begin{tabularx}{\linewidth}{@{}
  >{\raggedright\arraybackslash}p{0.14\linewidth}
  >{\raggedright\arraybackslash}p{0.32\linewidth}
  >{\raggedright\arraybackslash}X@{}}
\toprule
\rowcolor{UTPPaleBlue!92}
\textbf{Platform} & \textbf{Recommended route} & \textbf{What you use in the labs} \\
\midrule
Linux & Native Linux & Terminal, Vim, GCC, Bash \\
Windows & WSL 2 with Ubuntu & Ubuntu/WSL terminal \\
macOS & macOS Terminal & Terminal and compiler \\
\bottomrule
\end{tabularx}
\end{mltable}

\section*{Windows students: install WSL}
WSL lets Windows run a Linux environment alongside Windows. Your Windows desktop remains unchanged; the \textbf{Ubuntu/WSL terminal} becomes the laboratory environment.

Open \textbf{PowerShell or Windows Terminal as Administrator}:
\begin{lstlisting}[style=osbash]
wsl --install
\end{lstlisting}
Restart if Windows asks. Open \textbf{Ubuntu}, create a Linux username and password, then verify from PowerShell:
\begin{lstlisting}[style=osbash]
wsl --status
wsl --list --verbose
\end{lstlisting}

\begin{mlfigure}
\begin{tikzpicture}[
  box/.style={draw=UTPBlue,fill=UTPPaleBlue,rounded corners=2mm,
              text width=4.25cm,minimum height=1.55cm,inner sep=3mm,
              align=center,font=\small\color{UTPNavy}},
  linux/.style={box,draw=UTPTeal,fill=UTPPaleTeal},
  a/.style={-{Stealth[length=2.2mm]},thick,draw=UTPMuted}
]
\node[box] (win) at (-4.0,0) {\textbf{Windows PowerShell}\\[2pt]
  install and check WSL\\\texttt{wsl --status}};
\node[linux] (wsl) at (4.0,0) {\textbf{Ubuntu / WSL shell}\\[2pt]
  run course commands\\\texttt{gcc}, \texttt{vim}, \texttt{ps}, \texttt{bash}};
\draw[a] (win) -- node[above=3pt,fill=white,inner sep=1.5pt,
  text width=2.6cm,align=center,font=\scriptsize\color{UTPMuted}]
  {launches the\\Linux environment} (wsl);
\end{tikzpicture}
\end{mlfigure}

\begin{pitfall}
Do not type Linux package, compiler, or lab commands into PowerShell. After WSL is installed, run the practical work inside Ubuntu/WSL unless the guide explicitly says \textbf{PowerShell}.
\end{pitfall}

\section*{Install the command-line tools}
On Ubuntu, Debian, or WSL:
\begin{lstlisting}[style=osbash]
sudo apt update
sudo apt install build-essential gdb make vim git openssh-client
\end{lstlisting}
Verify the core tools:
\begin{lstlisting}[style=osbash]
gcc --version
vim --version | head -n 1
git --version
uname -a
\end{lstlisting}

For WSL, keep Linux course work in the Linux home directory rather than under \texttt{/mnt/c}:
\begin{lstlisting}[style=osbash]
mkdir -p ~/operating-systems
cd ~/operating-systems
pwd
\end{lstlisting}

\section*{The CLI commands you need first}
\begin{mltable}
\renewcommand{\arraystretch}{1.18}
\begin{tabularx}{\linewidth}{@{}
  >{\raggedright\arraybackslash}p{0.16\linewidth}
  >{\raggedright\arraybackslash}p{0.28\linewidth}
  >{\raggedright\arraybackslash}X@{}}
\toprule
\rowcolor{UTPPaleBlue!92}
\textbf{Command} & \textbf{Purpose} & \textbf{Example} \\
\midrule
\texttt{pwd} & show the current directory & \texttt{pwd} \\
\texttt{ls} & list files & \texttt{ls -la} \\
\texttt{cd} & change directory & \texttt{cd all\_experiments/Lab01} \\
\texttt{mkdir} & create a directory & \texttt{mkdir practice} \\
\texttt{cp} / \texttt{mv} & copy / move or rename & \texttt{cp a.c b.c}, \texttt{mv old new} \\
\texttt{rm} & remove a file & \texttt{rm program} \\
\texttt{man} & open the manual & \texttt{man ps} \\
\bottomrule
\end{tabularx}
\end{mltable}

\section*{Vim: the editor used in this course}
Vim is a terminal editor. You do not need advanced Vim knowledge; the course uses a small, repeatable set of commands.

\begin{mlfigure}
\begin{tikzpicture}[
  n/.style={draw=UTPBlue,fill=UTPPaleBlue,rounded corners=2mm,
            text width=4.25cm,minimum height=1.55cm,inner sep=3.5mm,
            align=center,font=\small\color{UTPNavy}},
  i/.style={n,draw=UTPTeal,fill=UTPPaleTeal,minimum height=2.05cm},
  c/.style={n,draw=UTPOrange,fill=UTPPaleOrange,minimum height=2.05cm},
  bileft/.style={{Stealth[length=2.5mm,width=2mm]}-{Stealth[length=2.5mm,width=2mm]},
                line width=1.05pt,draw=UTPTeal!85!black,
                shorten <=2pt,shorten >=2pt},
  biright/.style={{Stealth[length=2.5mm,width=2mm]}-{Stealth[length=2.5mm,width=2mm]},
                 line width=1.05pt,draw=UTPOrange!90!black,
                 shorten <=2pt,shorten >=2pt}
]
\node[n] (normal) at (0,2.65) {\textbf{NORMAL MODE}\\[3pt]
  navigate, edit, search, and issue commands};
\node[i] (insert) at (-4.25,-0.75) {\textbf{INSERT MODE}\\[3pt]
  type and edit program text\\[5pt]
  {\scriptsize\textbf{Enter:} \texttt{i}\quad\textbf{Return:} \texttt{Esc}}};
\node[c] (command) at (4.25,-0.75) {\textbf{COMMAND-LINE MODE}\\[3pt]
  save, quit, and change settings\\[5pt]
  {\scriptsize\textbf{Enter:} \texttt{:}\quad\textbf{Return:} \texttt{Enter/Esc}}};

\draw[bileft] (normal.south west)
  to[out=235,in=55,looseness=1.08] (insert.north east);
\draw[biright] (normal.south east)
  to[out=305,in=125,looseness=1.08] (command.north west);
\end{tikzpicture}
\end{mlfigure}

Open a file from the shell:
\begin{lstlisting}[style=osbash]
vim hello.c
\end{lstlisting}

\begin{mltable}
\renewcommand{\arraystretch}{1.18}
\begin{tabularx}{\linewidth}{@{}
  >{\raggedright\arraybackslash}p{0.22\linewidth}
  >{\raggedright\arraybackslash}X@{}}
\toprule
\rowcolor{UTPPaleBlue!92}
\textbf{Vim command} & \textbf{What it does} \\
\midrule
\texttt{i} & enter Insert mode and start typing \\
\texttt{Esc} & return to Normal mode \\
\texttt{:w} & save \\
\texttt{:q} & quit if nothing changed \\
\texttt{:wq} & save and quit \\
\texttt{:q!} & quit and discard unsaved changes \\
\texttt{u} / \texttt{Ctrl+r} & undo / redo \\
\texttt{/word} & search forward for \texttt{word}; press \texttt{n} for the next match \\
\texttt{:set number} & show line numbers \\
\texttt{:25} & jump to line 25, useful for compiler errors \\
\bottomrule
\end{tabularx}
\end{mltable}

For guided practice, run \texttt{vimtutor} once before Lab 1.

\section*{The edit-compile-run cycle}
Create a first program with \texttt{vim hello.c}:
\begin{lstlisting}[style=osc]
#include <stdio.h>

int main(void) {
    printf("My OS lab environment is ready.\n");
    return 0;
}
\end{lstlisting}
Compile and run:
\begin{lstlisting}[style=osbash]
gcc -Wall -Wextra -std=c11 hello.c -o hello
./hello
\end{lstlisting}

\begin{mlfigure}
\begin{tikzpicture}[
  flowcard/.style={draw=UTPBlue!65,fill=white,rounded corners=1.5mm,
                   text width=2.45cm,minimum height=1.35cm,
                   inner xsep=2mm,inner ysep=2.5mm,align=center,
                   font=\small\color{UTPNavy}},
  forward/.style={-{Stealth[length=2.3mm,width=1.8mm]},line width=0.9pt,
                  draw=UTPTeal,shorten <=2pt,shorten >=2pt},
  feedback/.style={-{Stealth[length=2.2mm,width=1.7mm]},semithick,dashed,
                   draw=UTPMuted,rounded corners=4pt,
                   shorten <=2pt,shorten >=2pt},
  feedbacklabel/.style={fill=white,inner xsep=2mm,inner ysep=1pt,
                       font=\scriptsize\color{UTPMuted}}
]
\node[flowcard] (edit) at (-5.2,0) {\textcolor{UTPBlue}{\textbf{01}}\\[2pt]
  \textbf{EDIT SOURCE}\\[3pt]\texttt{vim file.c}};
\node[flowcard] (save) at (-1.75,0) {\textcolor{UTPBlue}{\textbf{02}}\\[2pt]
  \textbf{SAVE}\\[3pt]\texttt{:wq}};
\node[flowcard] (build) at (1.75,0) {\textcolor{UTPBlue}{\textbf{03}}\\[2pt]
  \textbf{COMPILE}\\[3pt]\texttt{gcc ...}};
\node[flowcard] (run) at (5.2,0) {\textcolor{UTPBlue}{\textbf{04}}\\[2pt]
  \textbf{RUN AND INSPECT}\\[3pt]\texttt{./program}};

\draw[forward] (edit.east) -- (save.west);
\draw[forward] (save.east) -- (build.west);
\draw[forward] (build.east) -- (run.west);

\draw[feedback] (run.south) -- ++(0,-1.0)
  -- node[feedbacklabel,below]{revise one thing, then repeat} ++(-10.4,0)
  -- (edit.south);
\end{tikzpicture}
\end{mlfigure}

\section*{Git and SSH setup for course work}
Git records changes to your source files. SSH lets Git authenticate to GitHub or another Git host without using an account password for every push.

\subsection*{1. Install Git and OpenSSH}
\begin{mltable}
\renewcommand{\arraystretch}{1.18}
\begin{tabularx}{\linewidth}{@{}
  >{\raggedright\arraybackslash}p{0.16\linewidth}
  >{\raggedright\arraybackslash}X@{}}
\toprule
\rowcolor{UTPPaleBlue!92}
\textbf{System} & \textbf{Command} \\
\midrule
Linux & \texttt{sudo apt update \&\& sudo apt install git openssh-client} \\
WSL & Run the same Linux command \textbf{inside Ubuntu/WSL}, not PowerShell \\
macOS & \texttt{brew install git} then \texttt{git --version} \\
\bottomrule
\end{tabularx}
\end{mltable}

Confirm:
\begin{lstlisting}[style=osbash]
git --version
ssh -V
\end{lstlisting}

\subsection*{2. Configure your commit identity}
Use the name and email that should appear on your course commits:
\begin{lstlisting}[style=osbash]
git config --global user.name "Your Name"
git config --global user.email "you@example.com"
git config --global init.defaultBranch main
git config --global --list
\end{lstlisting}

\subsection*{3. Create an SSH key}
First check whether a key already exists:
\begin{lstlisting}[style=osbash]
ls -al ~/.ssh
\end{lstlisting}
If \texttt{id\_ed25519} and \texttt{id\_ed25519.pub} are absent, create them:
\begin{lstlisting}[style=osbash]
ssh-keygen -t ed25519 -C "you@example.com"
\end{lstlisting}
Press Enter to accept the default location. A passphrase is recommended.

Load the key:
\begin{lstlisting}[style=osbash]
eval "$(ssh-agent -s)"
ssh-add ~/.ssh/id_ed25519
\end{lstlisting}
On macOS, \texttt{ssh-add --apple-use-keychain \string~/.ssh/id\_ed25519} may be used instead.

\subsection*{4. Copy only the public key}
\begin{mltable}
\renewcommand{\arraystretch}{1.18}
\begin{tabularx}{\linewidth}{@{}
  >{\raggedright\arraybackslash}p{0.16\linewidth}
  >{\raggedright\arraybackslash}X@{}}
\toprule
\rowcolor{UTPPaleBlue!92}
\textbf{System} & \textbf{Copy command} \\
\midrule
Linux & \texttt{cat \string~/.ssh/id\_ed25519.pub} and copy the complete printed line \\
WSL & \texttt{clip.exe < \string~/.ssh/id\_ed25519.pub} \\
macOS & \texttt{pbcopy < \string~/.ssh/id\_ed25519.pub} \\
\bottomrule
\end{tabularx}
\end{mltable}

In GitHub, open \textbf{Settings $>$ SSH and GPG keys $>$ New SSH key}, give the key a recognizable name, and paste the public key.

\begin{pitfall}
Never share or commit \texttt{\string~/.ssh/id\_ed25519}. That is the \textbf{private key}. Only the file ending in \texttt{.pub} is copied to the Git hosting service.
\end{pitfall}

\subsection*{5. Test the SSH connection}
For GitHub:
\begin{lstlisting}[style=osbash]
ssh -T git@github.com
\end{lstlisting}
The first connection may ask you to confirm the host fingerprint. For GitLab, use \texttt{git@gitlab.com}.

\subsection*{6. Connect the course folder}
For a new empty GitHub repository, run these commands from the course root and replace \texttt{USERNAME/REPOSITORY}:
\begin{lstlisting}[style=osbash]
git init -b main
git add .
git status
git commit -m "Start operating-systems course work"
git remote add origin git@github.com:USERNAME/REPOSITORY.git
git push -u origin main
\end{lstlisting}

For later work:
\begin{lstlisting}[style=osbash]
git status
git add .
git commit -m "Complete Lab 1 process observations"
git push
\end{lstlisting}


